\section*{ОБОЗНАЧЕНИЯ И СОКРАЩЕНИЯ}

БД -- база данных.

ИС -- информационная система.

ИТ -- информационные технологии. 

КТС -- комплекс технических средств.

ОМТС -- отдел материально-технического снабжения. 

ПО -- программное обеспечение.

РП -- рабочий проект.

СУБД -- система управления базами данных.

ТЗ -- техническое задание.

ТП -- технический проект.

API (Application Programming Interface) -- интерфейс, через который клиент (браузер) обменивается JSON-данными с сервером по протоколу HTTP.

UI (User Interface) – Пользовательский Интерфейс. Это то, как выглядит продукт и как пользователь с ним взаимодействует. Включает в себя визуальные элементы (кнопки, иконки, текст), а также их расположение и эстетику.

UX (User Experience) – Пользовательский Опыт. Это все ощущения и впечатления пользователя, возникающие при взаимодействии с продуктом. UX охватывает удобство, эффективность, доступность и общее удовлетворение от использования.

UX/UI Design – Дизайн Пользовательского Опыта и Пользовательского Интерфейса. Комплексный подход к созданию продуктов, учитывающий как эстетику и функциональность, так и удобство и приятность использования.

DM (Direct Message) – Прямое Сообщение. Личная переписка между двумя пользователями.

PM (Private Message) – Личное Сообщение. То же, что и DM, используется взаимозаменяемо.

Wireframe – Каркасный Прототип. Схематичное представление интерфейса без деталей дизайна, фокусирующееся на расположении элементов и функциональности.

Mockup – Макет. Статическое изображение интерфейса с проработанным дизайном (цвета, шрифты, иконки), но без интерактивности.

Prototype – Прототип. Интерактивная модель интерфейса, которая позволяет имитировать взаимодействие пользователя с продуктом. Может быть низкой или высокой детализации.

GUI (Graphical User Interface) – Графический Пользовательский Интерфейс. Тип пользовательского интерфейса, который позволяет взаимодействовать с электронными устройствами с помощью графических иконок и визуальных индикаторов.